\section{Method}
\label{sec:method}

\subsection{Notation and probe}
\label{sec:method:notation}

Let $f : \mathcal{X} \to \mathbb{R}^{d}$ be a frozen vision encoder with
output dimension $d$, and define the L2-normalised representation
\begin{equation}
  z(x) \;=\; \frac{f(x)}{\lVert f(x) \rVert_{2}} \;\in\; \mathbb{S}^{d-1}.
\end{equation}
A \emph{stable} pair $(x_{c}, x_{t})$ shares a single underlying lesion image
$u$ subjected to two independent draws of a synthetic nuisance augmentation
$\nu \sim \mathcal{N}_{F}$ from family $F$:
\begin{equation}
  x_{c} = \nu_{c}(u), \qquad x_{t} = \nu_{t}(u), \qquad \nu_{c}, \nu_{t} \sim \mathcal{N}_{F}.
\end{equation}
A \emph{changing} pair pairs two distinct lesions $u_{1}, u_{2}$ matched on
diagnosis $\mathrm{dx}(u_{1}) = \mathrm{dx}(u_{2})$ and anatomical site
$\mathrm{site}(u_{1}) = \mathrm{site}(u_{2})$; each side receives an
independent nuisance draw. The predictor $g_{\theta} : \mathbb{R}^{d} \to
\mathbb{R}^{d}$ minimises the L2-MSE objective on stable training pairs:
\begin{equation}
  \mathcal{L}(\theta) \;=\; \frac{1}{|\mathcal{D}_{\mathrm{train}}^{\mathrm{stable}}|}\!\!
  \sum_{(x_{c}, x_{t}) \in \mathcal{D}_{\mathrm{train}}^{\mathrm{stable}}}
  \lVert g_{\theta}(z(x_{c})) - z(x_{t}) \rVert_{2}^{2}
  \;+\; \mathcal{R}(\theta),
  \label{eq:loss}
\end{equation}
where $\mathcal{R}$ is a predictor-specific regulariser. For the linear
predictor we use a soft \texttt{weight $-$ I} penalty,
$\mathcal{R}(\theta) = \lambda \, \lVert W - I \rVert_{F}^{2}$, which, in
combination with the identity warm-start ($W \approx I$ at step zero,
\S\ref{sec:method:predictor}), holds the function near the identity at
initialisation. For the 2-layer identity-residual MLP, the residual
weight $W_{2}$ is zero-initialised so the predictor is identity at step
zero by construction, and $\mathcal{R}$ is standard $L^{2}$ weight decay
on the predictor weights (decoupled, in the Adam configuration);
applying a $\lVert W_{2} - I \rVert_{F}^{2}$ penalty here would
contradict the zero-init identity warm-start. Per-variant values of
$\lambda$ are listed in Appendix~\ref{app:hyperparameters}. At
evaluation time, the score for a
pair is the cosine \emph{distance} between the L2-normalised prediction and
the L2-normalised target:
\begin{equation}
  s(x_{c}, x_{t}) \;=\; 1 \;-\; \langle \tilde{g}_{\theta}(z(x_{c})), z(x_{t}) \rangle, \qquad
  \tilde{g}_{\theta}(\cdot) := \frac{g_{\theta}(\cdot)}{\lVert g_{\theta}(\cdot) \rVert_{2}}.
  \label{eq:score}
\end{equation}
A correctly oriented predictor produces $s(x_{c}, x_{t}) < s(x_{c}', x_{t}')$
whenever $(x_{c}, x_{t})$ is stable and $(x_{c}', x_{t}')$ is changing;
$\auroc$ quantifies this ranking.

\subsection{Predictor variants}
\label{sec:method:predictor}

\paragraph{Linear with identity warm-start.}
$g_{\theta}(z) = W z + b$ with $W$ initialised to $I + \epsilon$,
$\epsilon_{ij} \overset{\text{iid}}{\sim} \mathcal{N}(0, 0.005^{2})$, and
$b = 0$. Under SGD with the soft \texttt{weight $-$ I} penalty above, the
predictor at step zero is identity-equivalent.

\paragraph{2-layer identity-residual MLP.}
$g_{\theta}(z) = z + W_{2} \,\mathrm{ReLU}(W_{1} z + b_{1}) + b_{2}$, with
$W_{2}$ zero-initialised so the residual term vanishes at step zero and
$g_{\theta} = \mathrm{identity}$. Hidden dimension $512$.

\subsection{Three nuisance families}
\label{sec:method:nuisance}

The augmentation suite specifies three deterministic synthetic nuisance
families, summarised in Table~\ref{tab:nuisance-families}. The disjointness
rule is operation-level: no transform type is shared between any two
families. The exact operation parameters are listed in
Appendix~\ref{app:nuisance-operations}.

\begin{table}[h]
\centering
\small
\begin{tabular}{@{}lp{0.65\linewidth}@{}}
\toprule
Family & Operation set \\
\midrule
\strongfam            & brightness $\times$ contrast $\times$ saturation, rotation $\pm 15^{\circ}$, scale $0.82$--$1.00$, translate $\pm 5\%$, horizontal flip, Gaussian blur, Gaussian noise, JPEG quality $45$--$70$ \\
\strongholdone        & hue shift $\pm 12^{\circ}$ (HSV rotation), posterise $4$--$6$ bits, unsharp-mask sharpen, motion blur (linear kernel $5$--$15$\,px), random rectangular erasing $6$--$18\%$ area, JPEG quality $20$--$40$ \\
\strongholdtwo        & gamma $0.6$--$1.6$, colour temperature $0.85$--$1.15$, radial vignette $25$--$55\%$ falloff, salt-and-pepper noise $0.5$--$2.5\%$, JPEG quality $80$--$95$ \\
\bottomrule
\end{tabular}
\caption{Three deterministic nuisance families, partitioned along the
operation axis. Train stable pairs draw from $\{\strongfam, \strongholdone\}$;
val and test stable pairs use $\strongholdtwo$.}
\label{tab:nuisance-families}
\end{table}

\subsection{Splits and pair construction}
\label{sec:method:splits}

We use lesion-ID-aware splits at fractions $0.70 / 0.15 / 0.15$, deterministic
from a fixed seed, with group counts $5{,}229 / 1{,}120 / 1{,}121$ across
HAM10000's $10{,}015$ images and $7{,}470$ unique lesions. Each split
contains $1{,}000$ stable plus $1{,}000$ changing pairs (total $6{,}000$
pairs per run). The strict same-diagnosis-site changing-pair policy requires
that a context lesion's changing partner be a distinct lesion with identical
diagnosis and identical anatomical site. Train stable pairs rotate
deterministically between \strongfam\ and \strongholdone\ by pair index; val
and test stable pairs use \strongholdtwo\ exclusively.

\subsection{Evaluation protocol}
\label{sec:method:eval}

We report $\auroc$ as the primary metric, with $1{,}000$-sample bootstrap
$95\%$ confidence intervals computed by per-pair resampling. AUPRC, equal-
error-rate threshold, and FPR at fixed TPR $= 0.8$ accompany every run, as do
three representation-health checks: prediction-norm mean and minimum,
dimension-variance mean and minimum, and a binary \texttt{collapsed} flag
triggered when either falls below $10^{-8}$. None of the nine reported runs
is flagged as collapsed.
